\begin{proof}[Proof of \cref{obj:thm:uniform-annotation-frontier}]
\begin{enumerate}
\item Choose numerical constants \(c_{\mathrm{lab}}>0\) and \(c_{\mathrm{aff}}>0\) from \cref{obj:thm:rare-label-floor,obj:lem:normalized-affine-testing}, respectively. For any public indices in the stated range, write
\[
N=n+m,\qquad
S=n\epsilon,\qquad
\ell=\log(\exp(1)+S),\qquad
x=\frac{d}{N\epsilon\ell}.
\]
Then \(N>0\), \(S>0\), and \(\ell>0\). By \cref{obj:def:rate-handle},
\[
r(n,m,d,\epsilon)=\min\{1,S^{-1}+x^2\},
\qquad
0<r(n,m,d,\epsilon)\le1.
\]
The two cited results supply
\[
c_{\mathrm{lab}}\,b(n,\epsilon)
\le R(n,m,d,\epsilon),
\qquad
b(n,\epsilon)=\min\{1,S^{-1}\},
\]
throughout the public index range, and
\[
c_{\mathrm{aff}}\min\{1,x^2\}
\le R(n,m,d,\epsilon)
\quad\text{when}\quad
S\ge\exp(4096),\qquad x^2>S^{-1}.
\]
Both constants are independent of the public indices. We next establish a uniform upper constant for the specified estimator.
% lean: CausalSmith.Stat.AnnotationRarearmFrontier.rare_label_floor
% lean: CausalSmith.Stat.AnnotationRarearmFrontier.normalized_affine_testing

\item First suppose \(S\ge\exp(4096)\), and use the pool lengths and tuning parameters from \cref{obj:def:estimator-handle}. Since \(\epsilon\le1/4\),
\[
12\le4\exp(4096)\le n.
\]
The integer splits satisfy
\[
h_0\ge\frac n3-1\ge\frac n4,\qquad
h_p\ge h_0,\qquad
h_p\le h_f\le h_p+1.
\]
Indeed, \(n-3h_0\in\{0,1,2\}\), and splitting the remaining records into two parts gives the last two inequalities. Consequently,
\[
M_p
=h_p+\lfloor m/2\rfloor
\ge\frac n4+\frac{m-1}{2}
\ge\frac N8,
\]
where the last inequality uses \(n\ge12\) and \(m\ge0\). Moreover,
\[
0\le M_f-M_p
=(h_f-h_p)+(m-2\lfloor m/2\rfloor)\le2.
\]
Thus
\[
h_0,M_p,M_f\ge1,\qquad
\frac n4\le h_0,\qquad
\frac N8\le M_p,M_f,\qquad
M_p\le M_f\le3M_p.
\]
For the prescribed intensities this yields
\[
u=\frac{h_0}{8}\ge\frac n{32},\qquad
t_p=\frac{M_p}{8}\ge\frac N{64},\qquad
t=\frac{M_f}{8}\ge\frac N{64},\qquad
\frac13\le\frac t{t_p}\le3.
\]
These bounds include \(m=0\).
% lean: CausalSmith.Stat.AnnotationRarearmFrontier.hybrid_finite_pool_sizes
% lean: CausalSmith.Stat.AnnotationRarearmFrontier.hybrid_finite_pool_intensities

\item Under \cref{obj:ass:data-product}, the three disjoint pools are independent. Their observation laws are \(P\), \(P_{XA}\), and \(P_{XA}\), respectively: projecting a complete record gives the same treatment--covariate law as an auxiliary record. Extend each pool by an independent iid continuation and introduce independent prefix lengths
\[
J_{\mathrm o}\sim\operatorname{Poi}(u),\qquad
J_{\mathrm p}\sim\operatorname{Poi}(t_p),\qquad
J_{\mathrm f}\sim\operatorname{Poi}(t),
\qquad
J_{\mathrm o},J_{\mathrm p},J_{\mathrm f}\in\mathbb N_0,
\]
independent of the observations.

Here is the count calculation for each pool. Let its finite alphabet be \(\mathsf Z_{\mathrm{pool}}\), its atom probabilities be \(\nu_{\mathrm{pool}}(z)\), and its prefix intensity be \(\lambda_{\mathrm{pool}}>0\), so that
\[
\nu_{\mathrm{pool}}(z)\ge0,\qquad
\sum_{z\in\mathsf Z_{\mathrm{pool}}}\nu_{\mathrm{pool}}(z)=1.
\]
For counts \(\kappa_z\in\mathbb N_0\), define
\[
\kappa_{\mathrm{tot}}
=\sum_{z\in\mathsf Z_{\mathrm{pool}}}\kappa_z.
\]
Conditioning on the prefix length and using the multinomial count formula gives
\[
\begin{aligned}
\Pr\{\text{every atom }z\text{ has count }\kappa_z\}
&=
e^{-\lambda_{\mathrm{pool}}}
\frac{\lambda_{\mathrm{pool}}^{\kappa_{\mathrm{tot}}}}
     {\kappa_{\mathrm{tot}}!}
\frac{\kappa_{\mathrm{tot}}!}{\prod_z\kappa_z!}
\prod_z\nu_{\mathrm{pool}}(z)^{\kappa_z}\\
&=
\prod_z
e^{-\lambda_{\mathrm{pool}}\nu_{\mathrm{pool}}(z)}
\frac{(\lambda_{\mathrm{pool}}\nu_{\mathrm{pool}}(z))^{\kappa_z}}
     {\kappa_z!}.
\end{aligned}
\]
We use \(0^0=1\); an atom of zero probability has count zero almost surely. This factorization, together with independence of the pools, shows that the prefix counts from \cref{obj:def:estimator-handle} are jointly independent and satisfy
\[
Z_{aj}\sim\operatorname{Poi}(uq_{aj}(P)),\qquad
J_{aj}\sim\operatorname{Poi}(t_ps_{aj}(P)),\qquad
K_{aj}\sim\operatorname{Poi}(ts_{aj}(P)).
\]
Thus \cref{obj:lem:uniform-hybrid-arm-risk} applies with precisely the prescribed \(L,B,k_0\).

Define, for \(a\in\{0,1\}\),
\[
\widehat\psi_a=\sum_{j=1}^dV_{aj},
\qquad
\psi_a=\sum_{j=1}^dp_j(P)\mu_{aj}(P),
\]
on these ideal prefixes. The cited result supplies a numerical constant \(C_{\mathrm{hyb}}>0\) such that
\[
\mathbb E[(\widehat\psi_a-\psi_a)^2]
\le
C_{\mathrm{hyb}}
\left\{S^{-1}+\left(\frac d{N\epsilon L}\right)^2\right\}.
\]
The binary conditional means and the covariate probabilities satisfy
\[
0\le\mu_{aj}(P)\le1,\qquad
p_j(P)\ge0,\qquad
\sum_{j=1}^dp_j(P)=1,
\]
including the zero convention for null cells. Hence
\[
-1
=\sum_jp_j(P)(-1)
\le\tau(P)
=\psi_1-\psi_0
\le\sum_jp_j(P)=1.
\]
Define the clipped ideal-prefix contrast by
\[
F_{\mathrm{Poi}}
=\operatorname{clip}_{[-1,1]}(\widehat\psi_1-\widehat\psi_0).
\]
Clipping toward an interval containing the target decreases squared loss, and
\[
\begin{aligned}
\mathbb E[(F_{\mathrm{Poi}}-\tau(P))^2]
&\le
\mathbb E[(\widehat\psi_1-\widehat\psi_0-\tau(P))^2]\\
&\le
2\mathbb E[(\widehat\psi_1-\psi_1)^2]
+2\mathbb E[(\widehat\psi_0-\psi_0)^2]\\
&\le
4C_{\mathrm{hyb}}
\left\{S^{-1}+\left(\frac d{N\epsilon L}\right)^2\right\}.
\end{aligned}
\]
The second inequality is pointwise, so it applies with the shared counts used by the two arm statistics.
% lean: CausalSmith.Stat.AnnotationRarearmFrontier.hybrid_ordered_prefix_arm_risk
% lean: CausalSmith.Stat.AnnotationRarearmFrontier.hybrid_ordered_prefix_contrast_risk

\item We calibrate the last bound to the numerical rate and bound the prefix overflow cost. From \(S\ge\exp(4096)\) and the definition of \(L\),
\[
L=\left\lfloor\frac{\log S}{1024}\right\rfloor\ge4,
\qquad
\log S<1024(L+1)\le1280L.
\]
Also \(S\ge2\) and \(\exp(1)\ge2\), which imply
\[
\exp(1)+S\le S\exp(1),
\qquad
\ell\le\log S+1\le1281L.
\]
All denominators below are positive, so
\[
0\le\frac d{N\epsilon L}\le1281x,
\qquad
S^{-1}+\left(\frac d{N\epsilon L}\right)^2
\le1281^2(S^{-1}+x^2).
\]

For every positive pool length \(h\), the inequality \(e^h\ge1+h\ge h\) gives
\[
e^{-h}\le h^{-1}.
\]
Applying this to the three pool lengths and using the bounds already established,
\[
\begin{aligned}
e^{-h_0}+e^{-M_p}+e^{-M_f}
&\le\frac4n+\frac8n+\frac8n\\
&=\frac{20}{n}
\le\frac5S,
\end{aligned}
\]
where the last inequality uses \(\epsilon\le1/4\). Define the numerical constant
\[
C_{\mathrm{large}}=4C_{\mathrm{hyb}}1281^2+5>0.
\]
Combining the ideal-prefix risk, degree comparison, and overflow bound yields
\[
\begin{aligned}
\mathbb E[(F_{\mathrm{Poi}}-\tau(P))^2]
+e^{-h_0}+e^{-M_p}+e^{-M_f}
&\le
4C_{\mathrm{hyb}}1281^2(S^{-1}+x^2)+5S^{-1}\\
&=
C_{\mathrm{large}}(S^{-1}+x^2)-5x^2\\
&\le C_{\mathrm{large}}(S^{-1}+x^2).
\end{aligned}
\]
% lean: CausalSmith.Stat.AnnotationRarearmFrontier.hybrid_finite_degree_rate
% lean: CausalSmith.Stat.AnnotationRarearmFrontier.hybrid_finite_overflow_rate
% lean: CausalSmith.Stat.AnnotationRarearmFrontier.hybrid_tuned_prefix_and_overflow_rate

\item Apply \cref{obj:lem:finite-prefix-transfer} to the three positive pool lengths \(h_0,M_p,M_f\), the clipped prefix statistic, and the target \(\tau(P)\in[-1,1]\). Its finite average is exactly the estimator in \cref{obj:def:estimator-handle}, since
\[
\Pr(J_{\mathrm o}=h',J_{\mathrm p}=h'',J_{\mathrm f}=h''')
=
e^{-u-t_p-t}
\frac{u^{h'}}{h'!}
\frac{t_p^{h''}}{h''!}
\frac{t^{h'''}}{h'''!}.
\]
The zero output on overflow agrees with the construction. Independence and the stated pool laws therefore give
\[
\begin{aligned}
\mathbb E_P[
(\widehat\tau(\mathcal D,U)-\tau(P))^2]
&\le
\mathbb E[(F_{\mathrm{Poi}}-\tau(P))^2]
+e^{-h_0}+e^{-M_p}+e^{-M_f}\\
&\le C_{\mathrm{large}}(S^{-1}+x^2).
\end{aligned}
\]
This proves the uncapped upper bound for every model law when \(S\ge\exp(4096)\).
% lean: CausalSmith.Stat.AnnotationRarearmFrontier.hybrid_large_information_risk

\item Define
\[
C_{\mathrm{up}}
=\max\{\max\{C_{\mathrm{large}},4\},\exp(4096)\}>0.
\]
For fixed public indices, the observed-data space is finite, so the deterministic estimator is Borel measurable. Its definition is independent of \(U\). In the averaging branch, every prefix statistic lies in \([-1,1]\), the displayed prefix probabilities are nonnegative, and
\[
\sum_{h'=0}^{h_0}\sum_{h''=0}^{M_p}\sum_{h'''=0}^{M_f}
\Pr(J_{\mathrm o}=h',J_{\mathrm p}=h'',J_{\mathrm f}=h''')
\le1.
\]
Thus, with the remaining probability assigned output zero,
\[
-1\le\widehat\tau(\mathcal D,U)\le1.
\]
The zero branch has the same range. Since \(\tau(P)\in[-1,1]\), this also gives the universal bound
\[
\mathbb E_P[
(\widehat\tau(\mathcal D,U)-\tau(P))^2]\le4.
\]

We now follow the three upper-bound regimes. If \(S<\exp(4096)\), the estimator is zero, and
\[
\mathbb E_P[
(\widehat\tau(\mathcal D,U)-\tau(P))^2]
=\tau(P)^2\le1.
\]
Moreover,
\[
e^{-4096}\le1,\qquad e^{-4096}\le S^{-1},
\qquad
e^{-4096}\le r(n,m,d,\epsilon),
\]
so
\[
\mathbb E_P[
(\widehat\tau(\mathcal D,U)-\tau(P))^2]
\le1
\le\exp(4096)r(n,m,d,\epsilon)
\le C_{\mathrm{up}}r(n,m,d,\epsilon).
\]
If \(S\ge\exp(4096)\) and \(1\le S^{-1}+x^2\), then
\[
r(n,m,d,\epsilon)=1,
\qquad
\mathbb E_P[
(\widehat\tau(\mathcal D,U)-\tau(P))^2]
\le4\le C_{\mathrm{up}}r(n,m,d,\epsilon).
\]
Finally, if \(S\ge\exp(4096)\) and \(S^{-1}+x^2<1\), the preceding uncapped bound gives
\[
\mathbb E_P[
(\widehat\tau(\mathcal D,U)-\tau(P))^2]
\le C_{\mathrm{large}}(S^{-1}+x^2)
\le C_{\mathrm{up}}r(n,m,d,\epsilon).
\]
Taking the supremum over the model class proves
\[
\sup_{P\in\mathcal M_{d,\epsilon}}
\mathbb E_P[
(\widehat\tau(\mathcal D,U)-\tau(P))^2]
\le C_{\mathrm{up}}r(n,m,d,\epsilon).
\]
% lean: CausalSmith.Stat.AnnotationRarearmFrontier.hybrid_upper

\item Choose the final numerical constants by
\[
E_{\mathrm{cut}}=\exp(4096),\qquad
c=\min\left\{
\frac{c_{\mathrm{lab}}}{2E_{\mathrm{cut}}},
\frac{c_{\mathrm{aff}}}{2}
\right\},
\qquad
C=\max\{C_{\mathrm{up}},c\}.
\]
They satisfy
\[
1\le E_{\mathrm{cut}},\qquad
0<c\le C<\infty,
\]
and the calibrations needed below are
\[
c\le\frac{c_{\mathrm{lab}}}{2E_{\mathrm{cut}}}
\le\frac{c_{\mathrm{lab}}}{E_{\mathrm{cut}}},
\qquad
2c\le\frac{c_{\mathrm{lab}}}{E_{\mathrm{cut}}}
\le c_{\mathrm{lab}},
\qquad
2c\le c_{\mathrm{aff}}.
\]
All these constants were chosen independently of the public indices.

The lifted rule \((\mathcal D,U)\mapsto\widehat\tau(\mathcal D,U)\) is Borel measurable by composition with the data projection and takes values in \([-1,1]\). It is therefore admissible in \cref{obj:def:risk}. Squared-error risks are nonnegative, and the defining minimax infimum yields
\[
R(n,m,d,\epsilon)
\le
\sup_{P\in\mathcal M_{d,\epsilon}}
\mathbb E_P[
(\widehat\tau(\mathcal D,U)-\tau(P))^2]
\le C_{\mathrm{up}}r(n,m,d,\epsilon)
\le Cr(n,m,d,\epsilon),
\]
where the last inequality uses \(r(n,m,d,\epsilon)\ge0\).
% lean: CausalSmith.Stat.AnnotationRarearmFrontier.uniform_annotation_frontier

\item For the lower bound, first suppose \(S<E_{\mathrm{cut}}\). Positivity of \(S\) and \(E_{\mathrm{cut}}\ge1\) gives
\[
E_{\mathrm{cut}}^{-1}\le S^{-1},
\qquad
E_{\mathrm{cut}}^{-1}\le1,
\qquad
E_{\mathrm{cut}}^{-1}
\le b(n,\epsilon).
\]
Using \(r(n,m,d,\epsilon)\le1\), the constant calibration, and the rare-label bound,
\[
\begin{aligned}
c\,r(n,m,d,\epsilon)
&\le c\\
&\le c_{\mathrm{lab}}E_{\mathrm{cut}}^{-1}\\
&\le c_{\mathrm{lab}}b(n,\epsilon)\\
&\le R(n,m,d,\epsilon).
\end{aligned}
\]
% lean: CausalSmith.Stat.AnnotationRarearmFrontier.uniform_annotation_frontier

\item Now suppose \(S\ge E_{\mathrm{cut}}\). Then \(S\ge1\), and hence
\[
b(n,\epsilon)=S^{-1}.
\]
If \(x^2\le S^{-1}\), the rate satisfies
\[
r(n,m,d,\epsilon)
\le S^{-1}+x^2\le2S^{-1}.
\]
It follows that
\[
c\,r(n,m,d,\epsilon)
\le2cS^{-1}
\le c_{\mathrm{lab}}S^{-1}
\le R(n,m,d,\epsilon).
\]

If \(x^2>S^{-1}\), all hypotheses of \cref{obj:lem:normalized-affine-testing} hold, and it supplies
\[
c_{\mathrm{aff}}\min\{1,x^2\}\le R(n,m,d,\epsilon).
\]
To compare this quantity with the rate, split according to whether \(x^2\le1\). In that case,
\[
r(n,m,d,\epsilon)
\le S^{-1}+x^2
<2x^2
=2\min\{1,x^2\}.
\]
When \(x^2>1\),
\[
r(n,m,d,\epsilon)\le1\le2=2\min\{1,x^2\}.
\]
Thus in both cases,
\[
\begin{aligned}
c\,r(n,m,d,\epsilon)
&\le2c\min\{1,x^2\}\\
&\le c_{\mathrm{aff}}\min\{1,x^2\}\\
&\le R(n,m,d,\epsilon).
\end{aligned}
\]
Together with the preceding branches, this proves the lower bound throughout the public index range.
% lean: CausalSmith.Stat.AnnotationRarearmFrontier.uniform_annotation_frontier

\item Finally, fix \(d\ge2\), \(0<\epsilon\le1/4\), and \(P\in\mathcal M_{d,\epsilon}\). Choose the independent conditional extension specified in \cref{obj:lem:causal-realization}: \(H\) is a joint law of \((X,A,Y,Y_0,Y_1)\), with binary potential outcomes, satisfying
\[
\mathcal L_H(X,A,Y)=P,\qquad
Y=AY_1+(1-A)Y_0\quad H\text{-almost surely},
\qquad
(Y_0,Y_1)\perp\!\!\!\perp A\mid X.
\]
The cited result supplies both this extension and, for every extension with the same observed marginal satisfying \cref{obj:ass:consistency,obj:ass:exchangeability}, the identity
\[
\mathbb E_H[Y_1-Y_0]=\tau(P).
\]
This establishes the causal realization and identification assertions.
% lean: CausalSmith.Stat.AnnotationRarearmFrontier.causal_realization
\end{enumerate}
\end{proof}
